
\documentclass[UTF8,zihao=-4]{ctexart}
\usepackage[a4paper,margin=2.15cm]{geometry}
\usepackage{amsmath,amssymb,bm,mathtools}
\usepackage{esint}
\usepackage{booktabs,longtable,tabularx,array,multicol}
\usepackage[most]{tcolorbox}
\usepackage{xcolor}
\usepackage{enumitem}
\usepackage{hyperref}
\usepackage{fancyhdr}
\usepackage{titlesec}
\usepackage{lastpage}
\usepackage{fvextra}
\usepackage{tikz}
\usepackage{eso-pic}
\usepackage{fontspec}
\setCJKmainfont{Noto Serif CJK SC}
\setCJKsansfont{Noto Sans CJK SC}
\setCJKmonofont{Noto Sans Mono CJK SC}
\setmainfont{DejaVu Serif}
\setsansfont{DejaVu Sans}
\setmonofont{DejaVu Sans Mono}
\hypersetup{colorlinks=true,linkcolor=blue!55!black,urlcolor=blue!55!black,citecolor=blue!55!black,pdftitle={高等数学A（二）考前复习知识点讲义},pdfauthor={Jerry}}
\definecolor{mainblue}{HTML}{1F4E79}
\definecolor{lightblue}{HTML}{EAF3FF}
\definecolor{warnred}{HTML}{B3261E}
\definecolor{lightred}{HTML}{FFF1F1}
\definecolor{darkgreen}{HTML}{1B6B3A}
\definecolor{lightgreen}{HTML}{EEF9F0}
\definecolor{orange}{HTML}{B26A00}
\definecolor{lightorange}{HTML}{FFF7E8}
\definecolor{softgray}{HTML}{777777}
\ctexset{section={format=\Large\bfseries\sffamily\color{mainblue}},subsection={format=\large\bfseries\sffamily\color{mainblue!85!black}},subsubsection={format=\normalsize\bfseries\sffamily\color{mainblue!75!black}}}
\titleformat{\paragraph}[runin]{\bfseries\sffamily\color{mainblue!75!black}}{}{0pt}{}[\quad]
\pagestyle{fancy}
\fancyhf{}
\lhead{高等数学A（二）考前复习知识点讲义}
\rhead{\leftmark}
\cfoot{\thepage/\pageref{LastPage}}
\setlength{\headheight}{15pt}
\setlist[itemize]{leftmargin=2em,itemsep=0.25em,topsep=0.25em}
\setlist[enumerate]{leftmargin=2em,itemsep=0.25em,topsep=0.25em}
\renewcommand{\arraystretch}{1.32}
\allowdisplaybreaks
\newcommand{\dd}{\,\mathrm d}
\newcommand{\ee}{\mathrm e}
\newcommand{\ii}{\mathrm i}
\newcommand{\ve}[1]{\bm{#1}}
\newcommand{\grad}{\nabla}
\newcommand{\epsz}{\varepsilon_0}
\newcommand{\muz}{\mu_0}
\newtcolorbox{keybox}[1][]{breakable,colback=lightblue,colframe=mainblue,arc=2mm,title=\textbf{#1},fonttitle=\bfseries\sffamily,coltitle=white}
\newtcolorbox{warnbox}[1][]{breakable,colback=lightred,colframe=warnred,arc=2mm,title=\textbf{#1},fonttitle=\bfseries\sffamily,coltitle=white}
\newtcolorbox{examplebox}[1][]{breakable,colback=lightgreen,colframe=darkgreen,arc=2mm,title=\textbf{#1},fonttitle=\bfseries\sffamily,coltitle=white}
\newtcolorbox{methodbox}[1][]{breakable,colback=lightorange,colframe=orange,arc=2mm,title=\textbf{#1},fonttitle=\bfseries\sffamily,coltitle=white}
\DefineVerbatimEnvironment{Code}{Verbatim}{fontsize=\small,breaklines,breakanywhere,frame=single,framesep=2mm}
\definecolor{watermarkgray}{HTML}{777777}
\AddToShipoutPictureBG{%
  \begin{tikzpicture}[remember picture,overlay]
    \node[rotate=35,scale=1.35,text=watermarkgray,opacity=0.14] at (current page.center) {Jerry 制作｜仅供个人复习使用｜禁止盗印与商用传播};
  \end{tikzpicture}%
}

\begin{document}
\begin{titlepage}
\centering
\vspace*{2.3cm}
{\Huge\bfseries\sffamily 高等数学A（二）考前复习知识点讲义\par}
\vspace{0.75cm}
{\Large 按两套模拟卷补配的考试导向版\par}
\vspace{0.45cm}
{\large 空间解析几何 - 多元函数微分法 - 重积分 - 曲线曲面积分 - 微分方程\par}
\vfill
\begin{tcolorbox}[colback=lightblue,colframe=mainblue,width=0.86\textwidth,arc=2mm]
\textbf{使用定位：}这份讲义按已出模拟卷的题型、分值与考点重新整理，重点放在高频结论、题型识别、固定步骤和易错点，不按教材逐段搬运。\par
\textbf{复习方法：}先看题型结构和优先级，再按模块刷小题和大题，最后用公式速查表与最后 30 分钟清单做查漏。
\end{tcolorbox}
\vfill
{\large 资料整理：Jerry\par}
{\large 版本：V1.1\par}
{\large \today\par}
\end{titlepage}
\tableofcontents
\newpage

\section{考试结构与复习策略}

\subsection{题型与分值结构}

\begin{longtable}{p{0.20\textwidth}p{0.16\textwidth}p{0.16\textwidth}p{0.14\textwidth}p{0.24\textwidth}}
\toprule
\textbf{题型} & \textbf{题量} & \textbf{单题分值} & \textbf{总分} & \textbf{主要考法} \\
\midrule
选择题 & 10 & 2 分 & 20 & 概念判断、公式识别、简单计算 \\
填空题 & 5 & 4 分 & 20 & 极限、偏导、Jacobi、面积体积、线积分短算 \\
计算/解答题 & 5 & 10 分 & 50 & 切线切平面、隐函数求导、重积分、线面积分、曲面积分 \\
证明题 & 1 & 10 分 & 10 & 梯度法向量、沿曲线常数、偏导条件证明 \\
\bottomrule
\end{longtable}

\begin{keybox}[拿分顺序]
先拿选择填空中的基础公式与判断题，再做计算题中固定套路题：隐函数求导、二重积分换序、三重积分坐标变换、Green/Gauss 公式。证明题最后做，至少写出参数化、求导或梯度正交的主线。
\end{keybox}

\subsection{考点优先级}

\begin{longtable}{p{0.16\textwidth}p{0.24\textwidth}p{0.26\textwidth}p{0.24\textwidth}}
\toprule
\textbf{优先级} & \textbf{模块} & \textbf{样卷出现方式} & \textbf{复习要求} \\
\midrule
必考核心 & 多元函数求导与隐函数 & 偏导、二阶偏导、全微分、切平面、梯度 & 会列方程、会隐求导、会检查可微条件 \\
必考核心 & 重积分 & 二重积分换序、极坐标、三重积分柱坐标/球坐标 & 会画区域、会选坐标、会用对称性 \\
必考核心 & 曲线/曲面积分 & Green 公式、路径无关、Gauss 公式、通量 & 会判定方向、边界和闭合性 \\
高频重点 & 空间解析几何 & 直线平面关系、曲线切线、曲面切平面 & 会用方向向量、法向量和梯度 \\
高频重点 & 微分方程 & 二阶常系数、重根、共振特解 & 会写特征方程与特解形式 \\
可能考点 & 曲面面积、旋转体体积、几何最值 & 填空或选择题 & 会套面积元、体积元和拉格朗日乘数法 \\
了解即可 & 理论性定义 & 连续、可微、路径无关的条件说明 & 能判断真伪，不深挖证明细节 \\
\bottomrule
\end{longtable}

\section{模块一：空间解析几何与切线切平面}

\subsection{核心公式}
\begin{keybox}[向量、直线、平面]
\begin{align*}
\text{平面：}\quad & A(x-x_0)+B(y-y_0)+C(z-z_0)=0,\quad \ve n=(A,B,C),\\
\text{直线：}\quad & \frac{x-x_0}{l}=\frac{y-y_0}{m}=\frac{z-z_0}{n},\quad \ve v=(l,m,n),\\
\text{点到平面距离：}\quad & d=\frac{|Ax_0+By_0+Cz_0+D|}{\sqrt{A^2+B^2+C^2}}.
\end{align*}
直线与平面的位置优先看 $\ve v\cdot\ve n$：不为 0 时相交；等于 0 时再代点判断平行或在平面内。
\end{keybox}

\begin{methodbox}[空间曲线切线与法平面]
若曲线 $C$ 由两个曲面 $F(x,y,z)=0, G(x,y,z)=0$ 相交而成，在点 $P$ 处
\[
\ve v=\grad F(P)\times \grad G(P).
\]
切线方程用点向式，法平面方程为
\[
\ve v\cdot[(x,y,z)-P]=0.
\]
\end{methodbox}

\begin{methodbox}[曲面切平面与法线]
显式曲面 $z=f(x,y)$ 在 $(x_0,y_0,z_0)$ 处：
\[
z-z_0=f_x(x_0,y_0)(x-x_0)+f_y(x_0,y_0)(y-y_0).
\]
隐式曲面 $F(x,y,z)=0$：
\[
\grad F(P)\cdot[(x,y,z)-P]=0,\qquad \frac{x-x_0}{F_x(P)}=\frac{y-y_0}{F_y(P)}=\frac{z-z_0}{F_z(P)}.
\]
\end{methodbox}

\begin{warnbox}[易错点]
\begin{enumerate}
  \item 法向量与方向向量不能混用；直线方向向量垂直于平面法向量才可能平行于平面。
  \item 曲面 $z=f(x,y)$ 的切平面常写成 $f_x(x-x_0)+f_y(y-y_0)-(z-z_0)=0$，符号别反。
  \item 两曲面交线的切向量是两个法向量的叉乘，不是点乘。
\end{enumerate}
\end{warnbox}

\section{模块二：多元函数极限、偏导、全微分与隐函数}

\subsection{极限连续与可微判断}
\begin{keybox}[判断链]
\[
\text{可微} \Rightarrow \text{连续} \Rightarrow \text{极限存在},\qquad \text{可微} \Rightarrow \text{偏导存在}.
\]
反过来一般不成立。判断二元极限时，先试直线、抛物线等路径；若要证明存在，常用夹逼、极坐标或估阶。
\end{keybox}

\begin{methodbox}[二元极限题模板]
\begin{enumerate}
  \item 先代直线 $y=kx$ 或特殊路径，看结果是否依赖 $k$。
  \item 若怀疑存在，改用 $r=\sqrt{x^2+y^2}$ 估阶。
  \item 遇到 $\frac{x^ay^b}{x^2+y^2}$，常用 $|xy|\leq \frac{x^2+y^2}{2}$。
\end{enumerate}
\end{methodbox}

\subsection{方向导数、梯度与极值}
\begin{keybox}[梯度]
\[
\grad f=(f_x,f_y,f_z),\qquad D_{\ve l}f=\grad f\cdot \ve l.
\]
函数增长最快方向是梯度方向，最大方向导数为 $|\grad f|$。约束最值常用拉格朗日乘数：
\[
\grad f=\lambda \grad g.
\]
\end{keybox}

\begin{keybox}[二元极值判别]
设 $f_x=f_y=0$，
\[
A=f_{xx},\quad B=f_{xy},\quad C=f_{yy},\quad \Delta=AC-B^2.
\]
$\Delta>0,A>0$ 为极小；$\Delta>0,A<0$ 为极大；$\Delta<0$ 非极值；$\Delta=0$ 需另判。
\end{keybox}

\subsection{隐函数求导}
\begin{methodbox}[一个方程确定 $z=z(x,y)$]
设 $F(x,y,z)=0$ 且 $F_z\ne0$，则
\[
z_x=-\frac{F_x}{F_z},\qquad z_y=-\frac{F_y}{F_z}.
\]
求二阶偏导时，不要把 $z$ 当常数。可先由 $F_x+F_z z_x=0$ 对 $x$ 再求导。
\end{methodbox}

\begin{examplebox}[二阶隐偏导套路]
若 $F(x,y,z)=0$，求 $z_{xx}$：
\[
F_x+F_z z_x=0.
\]
对 $x$ 求导得
\[
F_{xx}+2F_{xz}z_x+F_{zz}z_x^2+F_z z_{xx}=0,
\]
因此
\[
z_{xx}=-\frac{F_{xx}+2F_{xz}z_x+F_{zz}z_x^2}{F_z}.
\]
\end{examplebox}

\section{模块三：二重积分与三重积分}

\subsection{二重积分}
\begin{methodbox}[二重积分固定步骤]
\begin{enumerate}
  \item 画出区域，先判断是 $x$ 型、$y$ 型还是适合极坐标。
  \item 检查奇偶性和区域对称性：奇函数在对称区域积分为 0。
  \item 若出现 $x^2+y^2$、圆盘、扇形，优先极坐标：$x=r\cos\theta,y=r\sin\theta,\dd x\dd y=r\dd r\dd\theta$。
  \item 换序题先把原积分翻译成不等式，再重新投影。
\end{enumerate}
\end{methodbox}

\begin{keybox}[常见区域]
\begin{itemize}
  \item $0\le x\le 1, x^2\le y\le x$：由 $y=x^2$ 与 $y=x$ 围成。
  \item 圆盘 $x^2+y^2\le R^2$：$0\le r\le R,0\le\theta\le2\pi$。
  \item 上半圆：$0\le r\le R,0\le\theta\le\pi$。
\end{itemize}
\end{keybox}

\subsection{三重积分}
\begin{keybox}[坐标变换]
柱坐标：
\[
x=r\cos\theta,\quad y=r\sin\theta,
\quad \dd V=r\dd r\dd\theta\dd z.
\]
球坐标：
\[
x=\rho\sin\varphi\cos\theta,
\quad y=\rho\sin\varphi\sin\theta,
\quad z=\rho\cos\varphi,
\quad \dd V=\rho^2\sin\varphi\dd\rho\dd\varphi\dd\theta.
\]
\end{keybox}

\begin{methodbox}[三重积分选坐标]
\begin{itemize}
  \item 球、半球、球壳：优先球坐标。
  \item 圆柱、旋转抛物面、$x^2+y^2$：优先柱坐标。
  \item 被积函数为奇函数且区域关于对应坐标平面对称：积分为 0。
\end{itemize}
\end{methodbox}

\section{模块四：曲线积分、曲面积分与积分公式}

\subsection{第一类与第二类曲线积分}
\begin{keybox}[曲线积分]
第一类曲线积分：
\[
\int_L f(x,y)\dd s.
\]
与方向无关，参数化后 $\dd s=\sqrt{(x'(t))^2+(y'(t))^2}\dd t$。
第二类曲线积分：
\[
\int_L P\dd x+Q\dd y.
\]
与方向有关，反向变号。
\end{keybox}

\begin{methodbox}[Green 公式]
若 $L$ 是平面区域 $D$ 的正向边界，
\[
\oint_L P\dd x+Q\dd y=\iint_D\left(\frac{\partial Q}{\partial x}-\frac{\partial P}{\partial y}\right)\dd x\dd y.
\]
正向边界指沿曲线前进时区域在左侧。若边界不是闭合曲线，不能直接套 Green 公式。
\end{methodbox}

\begin{methodbox}[路径无关]
在单连通区域内，若 $P_y=Q_x$，则 $P\dd x+Q\dd y$ 为全微分，积分只与端点有关。找势函数 $u$：先由 $u_x=P$ 对 $x$ 积分，再用 $u_y=Q$ 确定剩余函数。
\end{methodbox}

\subsection{曲面积分与 Gauss 公式}
\begin{keybox}[通量积分]
\[
\iint_{\Sigma} P\dd y\dd z+Q\dd z\dd x+R\dd x\dd y=\iint_{\Sigma}\ve F\cdot\ve n\dd S.
\]
闭曲面外侧用 Gauss 公式：
\[
\iint_{\partial\Omega}\ve F\cdot\ve n\dd S=\iiint_{\Omega}\nabla\cdot\ve F\dd V.
\]
\end{keybox}

\begin{methodbox}[曲面面积]
显式曲面 $z=f(x,y)$ 在 $D$ 上方：
\[
S=\iint_D\sqrt{1+f_x^2+f_y^2}\dd x\dd y.
\]
若题目给平面 $z=ax+by+c$，面积因子为 $\sqrt{1+a^2+b^2}$。
\end{methodbox}

\section{模块五：常微分方程与证明题}

\subsection{二阶常系数线性微分方程}
\begin{methodbox}[齐次方程]
对 $y''+py'+qy=0$ 写特征方程 $r^2+pr+q=0$：
\begin{itemize}
  \item 两个不同实根：$y=C_1e^{r_1x}+C_2e^{r_2x}$；
  \item 重根 $r$：$y=(C_1+C_2x)e^{rx}$；
  \item 共轭复根 $\alpha\pm \beta i$：$y=e^{\alpha x}(C_1\cos\beta x+C_2\sin\beta x)$。
\end{itemize}
\end{methodbox}

\begin{warnbox}[非齐次特解易错]
右端 $e^{\lambda x}$、$\sin bx$、$\cos bx$ 与齐次解重复时要乘 $x$；重复阶数越高，乘的 $x$ 次数越高。
\end{warnbox}

\subsection{证明题常见路线}
\begin{methodbox}[沿曲线常数证明]
若要证明 $f$ 在曲线 $\gamma(t)$ 上为常数，只需证明
\[
\frac{\dd}{\dd t}f(\gamma(t))=\grad f(\gamma(t))\cdot\gamma'(t)=0.
\]
圆周 $x=r\cos t,y=r\sin t$ 上，$\gamma'(t)=(-r\sin t,r\cos t)$。
\end{methodbox}

\begin{methodbox}[梯度垂直切向量]
曲面 $F(x,y,z)=0$ 上取任意曲线 $\ve r(t)$，有 $F(\ve r(t))\equiv0$。两边求导：
\[
\grad F(P)\cdot \ve r'(t_0)=0.
\]
故 $\grad F(P)$ 垂直于任意切向量，是曲面的法向量。
\end{methodbox}

\section{高频题型方法模板}

\begin{methodbox}[题型一：换序与换坐标]
\textbf{识别特征：}积分限中出现 $x^2,y/2,\sqrt y$ 等，或被积函数直接积分困难。\par
\textbf{固定步骤：}写不等式，画区域，投影到另一坐标轴，重写上下限；圆域优先极坐标。\par
\textbf{易错点：}上下限交点漏算；换极坐标忘记 Jacobian $r$。
\end{methodbox}

\begin{methodbox}[题型二：曲线/曲面积分公式选择]
\textbf{识别特征：}闭合平面曲线用 Green；闭曲面通量用 Gauss；路径无关看 $P_y=Q_x$。\par
\textbf{固定步骤：}先判断闭合性和方向，再写公式，再计算区域积分。\par
\textbf{易错点：}边界反向会变号；曲面不是闭合时不能直接 Gauss，需补面。
\end{methodbox}

\begin{methodbox}[题型三：隐函数二阶偏导]
\textbf{识别特征：}题干给 $F(x,y,z)=0$ 并要求 $z_x,z_y,z_{xx},z_{xy}$。\par
\textbf{固定步骤：}先确认 $F_z\ne0$；写一阶公式；再对一阶方程求导；最后代点。\par
\textbf{易错点：}二阶求导时 $z_x$ 也是函数，不能丢链式项。
\end{methodbox}

\section{公式速查表}

\begin{longtable}{p{0.23\textwidth}p{0.34\textwidth}p{0.33\textwidth}}
\toprule
\textbf{模块} & \textbf{公式/结论} & \textbf{使用场景} \\
\midrule
空间几何 & $d=\frac{|Ax_0+By_0+Cz_0+D|}{\sqrt{A^2+B^2+C^2}}$ & 点到平面距离 \\
曲面切平面 & $\grad F(P)\cdot[(x,y,z)-P]=0$ & 隐式曲面 \\
梯度 & $D_{\ve l}f=\grad f\cdot\ve l$ & 方向导数、最快增长 \\
隐函数 & $z_x=-F_x/F_z,z_y=-F_y/F_z$ & 一个方程确定 $z$ \\
二重积分 & $\dd x\dd y=r\dd r\dd\theta$ & 极坐标换元 \\
三重积分 & $\dd V=r\dd r\dd\theta\dd z$ & 柱坐标 \\
Green & $\oint P\dd x+Q\dd y=\iint(Q_x-P_y)\dd A$ & 平面闭曲线正向边界 \\
Gauss & $\iint_{\partial\Omega}\ve F\cdot\ve n\dd S=\iiint_{\Omega}\nabla\cdot\ve F\dd V$ & 闭曲面外侧通量 \\
曲面面积 & $S=\iint_D\sqrt{1+f_x^2+f_y^2}\dd A$ & $z=f(x,y)$ \\
ODE & $r^2+pr+q=0$ & 二阶常系数齐次方程 \\
\bottomrule
\end{longtable}

\section{易错点清单}
\begin{warnbox}[高频扣分点]
\begin{enumerate}
  \item 连续、偏导存在、可微三者关系方向写反。
  \item 二重积分换序时只抄上下限，不画区域。
  \item 极坐标、柱坐标、球坐标漏乘 Jacobian。
  \item Green 公式边界方向写错，正向是逆时针。
  \item 曲面积分方向不看外侧/上侧/下侧。
  \item 路径无关只看 $P_y=Q_x$，忘记区域要单连通。
  \item 隐函数二阶偏导漏掉 $F_{xz}z_x$、$F_{zz}z_x^2$ 等链式项。
\end{enumerate}
\end{warnbox}

\section{考前 3～7 天复习计划}
\begin{methodbox}[7 天版]
\begin{enumerate}
  \item Day 1：空间几何、切线切平面、梯度方向导数。
  \item Day 2：多元极限、偏导、全微分、隐函数求导。
  \item Day 3：二重积分换序、极坐标和对称性。
  \item Day 4：三重积分、曲面面积和几何应用。
  \item Day 5：曲线积分、Green 公式、路径无关。
  \item Day 6：曲面积分、Gauss 公式、微分方程。
  \item Day 7：完整做一套模拟卷，整理错题并背公式速查表。
\end{enumerate}
\end{methodbox}

\begin{methodbox}[3 天版]
\begin{enumerate}
  \item Day 1：背公式速查表，刷选择填空。
  \item Day 2：集中做隐函数、重积分、线面积分大题。
  \item Day 3：复盘错题，练证明题模板，背最后 30 分钟清单。
\end{enumerate}
\end{methodbox}

\section{考前最后 30 分钟速记清单}
\begin{keybox}[只看这些]
\begin{enumerate}
  \item 切平面：显式 $z-z_0=f_x(x-x_0)+f_y(y-y_0)$；隐式 $\grad F(P)$ 是法向量。
  \item 隐函数：$z_x=-F_x/F_z,z_y=-F_y/F_z$。
  \item 极坐标 $\dd A=r\dd r\dd\theta$；柱坐标 $\dd V=r\dd r\dd\theta\dd z$；球坐标 $\dd V=\rho^2\sin\varphi\dd\rho\dd\varphi\dd\theta$。
  \item Green：$Q_x-P_y$；Gauss：散度 $P_x+Q_y+R_z$。
  \item 路径无关：$P_y=Q_x$ 加单连通；积分等于势函数端点差。
  \item ODE 重根用 $(C_1+C_2x)e^{rx}$；共振特解要乘 $x$。
  \item 大题先写公式和区域，再算；方向、符号、Jacobian 是最后检查项。
\end{enumerate}
\end{keybox}

\end{document}
